\chapter*{РЕФЕРАТ}
\addcontentsline{toc}{chapter}{РЕФЕРАТ}

Расчетно-пояснительная записка \begin{NoHyper}\pageref{LastPage}\end{NoHyper}~с., \totfig~рис., \tottab~табл., 1~прил., \mbox{\thetotbibentries~источн.}

ДОСУГ, МЕРОПРИЯТИЕ, БАЗА ДАННЫХ, РЕЛЯЦИОННАЯ МОДЕЛЬ, СИСТЕМА УПРАВЛЕНИЯ БАЗАМИ ДАННЫХ, POSTGRESQL, ТРИГГЕР.

Цель работы -- спроектировать и реализовать базу данных для поиска и организации досуга и приложение для взаимодействия с ней.

В аналитической части рассмотрены задачи поиска мероприятий и организации совместного досуга, выполнено сравнение известных решений. Сформулированы функциональные и нефункциональные требования, выделены сущности и связи между ними, определены роли пользователей. По результатам сравнения моделей данных обоснован выбор реляционной модели.

В конструкторской части разработана структура БД для хранения пользователей, мероприятий, мест проведения, категорий, изображений, участия и отзывов. Определены ограничения целостности, проверено соответствие третьей нормальной форме. Разработаны ролевая модель, алгоритм триггера сохранения фактов удаления объявлений и архитектура приложения.

В технологической части обоснован выбор PostgreSQL и средств разработки приложения JoinUp. Реализованы серверная и клиентская части, ограничения БД, разграничение доступа и триггер. Описаны методы тестирования и результаты проверки создания мероприятия, отдельных операций приложения, ограничений целостности, прав БД и триггерной функции.

В исследовательской части определена зависимость времени выполнения запроса поиска по городу и диапазону дат от количества мероприятий с индексом и без него. Для 10000~мероприятий среднее время составило 4,905~мс без индекса и 3,038~мс с ним. В условиях проведенного исследования индекс сократил время выполнения запроса примерно на 38\%.

